\subsection*{13.3 Properties of Conditional Expectation}

In this section, we consider a probability space .(\Omega, \mathcal{F},P) and a sub-\sigma-field. \mathcal{G}of. \mathcal{F}.

We first address the existence of conditional expectation, which follows from the

Radon-Nikodym theorem.

Let \mu and \nu be measures defined on the same measurable space. We say that

measure \nu is absolutely continuous with respect to measure \mu if \mu(A)= 0

\nu(A)= 0 This is denoted by \nu\ll \mu .

For example, a continuous-type distribution is absolutely continuous with respect

to the Lebesgue measure. The theorem by Radon and Nikodym provides a partial

converse result. We state it below without proof.

\begin{thmbox}{13.5 (Radon-Nikodym)}
\end{thmbox}

Let. \mathcal{G}be a sub-\sigma-field of. \mathcal{F}Suppose. \nu and. \mu are\sigma-finite measures such that

\nu\ll \mu Then there exists a nonnegative \mathcal{G}-measurable function f such that

\nu(B)=

\int

B fd \mu for all B\in \mathcal{G}The function f is unique in the sense that if

g is another\mathcal{G}-measurable function with the same property, then g is equal to

f \mu-almost everywhere.

The function f is called the density or the Radon-Nikodym derivative of

\nu with respect to. \mu. It is represented by the symbol d\nu

d\mu .

\begin{thmbox}{13.6 (Existence of Conditional Expectation)}
\end{thmbox}

If X\in L 1(P) and \mathcal{G}is a sub-\sigma-field of \mathcal{F}, then the conditional expectation

E[X\vert\mathcal{G}] exists.

\textit{Proof} For random variable X in L2(P) , we know that the conditional expectation

of X can be obtained by the projection function.

For random variable X in L1(P) , the existence of conditional expectation is a

consequence of the Radon-Nikodym theorem because the measure \nu defined by

\nu(B)=

\int

B

XdP

is absolutely continuous with respect to measure P We can take the Radon-

Nikodym derivatived\nu/dP as the conditional expectation. \hfill $\square$

Most of the random variables we encounter have finite second moment. The

construction of conditional expectation using L2 theory and projection suffices for

most purposes. However, when a random variable is in L1 but not in. L2, we need to

use the Radon-Nikodym theorem to prove the existence. We refer readers to more

advanced textbooks, such as [ 2] and [4], for further details.

From now on, we will assume that conditional expectation for random variable

in. L1 exists.

\begin{thmbox}{13.7 (Properties of Conditional Expectation)}
\end{thmbox}

Let.(\Omega, \mathcal{F},P) denote a probability space, and let \mathcal{G}be a\sigma-subfield of. \mathcal{F}:

1. If X is inL1(P) and is\mathcal{G}-measurable, thenE[X\vert\mathcal{G}]= X as. In particular ,

for any constant c,E[c\vert\mathcal{G}]= c as.

2. (Nonnegativity) IfX \geq 0 as., thenE[X\vert\mathcal{G}]\geq 0 as.

3. (Linearity) SupposeX, Y \in L1(P) and a is a constant. We have

E[X + Y\vert\mathcal{G}]= E[X\vert\mathcal{G}]+ E[Y\vert\mathcal{G}],

and

E[aX\vert\mathcal{G}]= aE[X\vert\mathcal{G}].

4. (Monotonicity) If X and Y are in L1(P) and X \leq Y as., then E[X\vert\mathcal{G}]\leq

E[Y\vert\mathcal{G}] as.

5. (Conditional triangle inequality) For random variable X inL1(P) ,

\vert\vert\vertE[X\vert\mathcal{G}]

\vert\vert\vert \leqE

[

\vertX\vert\vert \mathcal{G}

]

\textit{Proof}

(1) We verify that X is \mathcal{G}-measurable, and we can replace X by X in the defining

propertyE[ X1B]= E[X1B] for allB \in\mathcal{G}The second statement in (1) follows

because a constant random variable is \mathcal{G}-measurable for any \sigma-algebra. \mathcal{G}.

(2) SupposeX \geq 0 almost surely, and. X is a version of E[X\vert\mathcal{G}]For integern \geq 1,

take En to be the set .{ : X \leq- 1/n}, for integer n \geq 1The s e t En is

\mathcal{G}-measurable because X is \mathcal{G}-measurable. Because X is nonnegative as. by

assumption,

\int

En

X=

\int

En

X \geq 0

for all n. Hence,

0 \leq

\int

En

XdP \leq

\int

En

(

- 1

n

)

dP=- 1

nP(En).

This is possible only if P(En) = 0. Since P(En) = 0 for all n \geq 1, the event

{ X< 0}=\cup nEn must have probability zero. This proves that X\geq 0 as.

(3) For any eventB\in \mathcal{G},we have

\int

B

X=

\int

B

X. (13.6)

\int

B

Y =

\int

B

Y, (13.7)

where X and Y are versions of E[X\vert\mathcal{G}] and E[Y\vert\mathcal{G}], respectively. By adding

(13.6) and (13.7), we obtain

\int

B

Z=

\int

B

(X+ Y) =

\int

B

( X+ Y). (13.8)

Comparing with the defining property of conditional expectation for E[Z\vert\mathcal{G}],

\int

B

Z=

\int

B

Z,

where. Z denotes a version of E[Z\vert\mathcal{G}], we see that

\int

B

Z=

\int

B

( X+ Y)

for allB\in \mathcal{G}This proves Z= X+ Y as.

Similarly, we can proveE[aX\vert\mathcal{G}]= a X as.

(4) SinceY- X \geq 0 as., by the nonnegativity of conditional expectation, we have

E[Y- X \vert\mathcal{G}]\geq 0 as. By linearity, we obtain E[Y\vert\mathcal{G}]- E[X\vert\mathcal{G}]\geq 0.

(5) SinceX\leq\vert X\vert, by monotonic property of conditional expectation,

E[X\vert\mathcal{G}]\leq E[\vert X\vert\vert \mathcal{G}]as.

Similarly, from.-X\leq\vert X\vert, we derive that

-E [X\vert\mathcal{G}]= E[-X\vert\mathcal{G}]\leq E[\vert X\vert\vert \mathcal{G}]as.

This proves that the absolute value of the conditional expectation of X given \mathcal{G}

is less than or equal to the conditional expectation of \vertX\vert given. \mathcal{G}.

\hfill $\square$

The properties listed in the previous theorem all have analog in ordinary

expectation. The basic convergence theorems such as Fatou's lemma, monotone

convergence theorem, and dominated convergence theorem also have versions for

conditional expectation. However, the next three properties are unique to conditional

expectation.

\begin{thmbox}{13.8 (Tower Property)}
\end{thmbox}

Suppose\mathcal{F}\mathcal{G} \mathcal{H} is a tower of \sigma-algebras and X is inL1(P) Then

E[X\vert\mathcal{H}]= E[E [X\vert\mathcal{G}]\vert\mathcal{H}]as.

\textit{Proof} We want to show that for all B\in \mathcal{H} ,

\int

B

E[X\vert\mathcal{H}]dP =

\int

B

E[E [X\vert\mathcal{G}]\vert\mathcal{H}]dP. (13.9)

The left-hand side is .

\int

B XdP The right-hand side is equal to .

\int

B E[X\vert\mathcal{G}]dP =\int

B XdP The first equality is due to B\in \mathcal{H} The second equality follows from

B\in \mathcal{G}Therefore, both sides of (13.9) are identical.

Since (13.9) holds for all \mathcal{H}-measurable set B , we conclude that E[E [X\vert\mathcal{G}]\vert\mathcal{H}]

is a version of the conditional expectation of X given. \mathcal{H}\hfill $\square$

\begin{thmbox}{13.9}
\end{thmbox}

Suppose XY are in L1(P,\mathcal{F}), where X is a \mathcal{G}-measurable function and Y is

\mathcal{F}-measurable and integrable, then

E[XY\vert\mathcal{G}]= XE[Y\vert\mathcal{G}]as.

In particular , when X is a function of Z , we have E[XY\vertZ]= XE[Y\vertZ] as.

\textit{Proof} We first show that, for any \mathcal{G}-measurable set A ,

\int

B

E[1AY\vert\mathcal{G}]dP =

\int

B

1AE[Y\vert\mathcal{G}]dP B\in \mathcal{G}. (13.10)

For eachB\in \mathcal{G}, the right-hand side of (13.10) equals

\int

1B\capAE[Y\vert\mathcal{G}]dP =

\int

B\capA

E[Y\vert\mathcal{G}]dP =

\int

B\capA

Yd P . ( B \cap A \in\mathcal{G}).

Meanwhile, the left-hand side of (13.10) is equal to

\int

B

1AYd P =

\int

1B\capAYd P =

\int

B\capA

Yd P . ( B \cap A \in\mathcal{G}).

This proves that (13.10) holds for allB\in \mathcal{G}, and hence for all \mathcal{G}-measurable simple

functions. Extend (13.10) to nonnegative functions and real-valued functions. \hfill $\square$

\begin{thmbox}{13.10}
\end{thmbox}

Suppose X is a random variable in .(\Omega,\mathcal{F},P) that is integrable and \mathcal{G}is a

sub-\sigma-algebra of. \mathcal{F}I f\sigma(X) and. \mathcal{G}are independent, then

E[X\vert\mathcal{G}]= E[X]as.

In particular , when X and Y are independent random variables, we have

E[X\vertY]= E[X] as.

\textit{Proof} It is sufficient to show that for each B\in \mathcal{G},

\int

B

XdP =

\int

B

E[X]dP. (13.11)

SupposeX= 1 A for some\mathcal{F}-measurable set A For any set B\in \mathcal{G}, the left-hand

side of (13.11) is

\int

B

1A dP=

\int

1A\capB dP= P(A \capB) = P(A)P(B).

On the other hand, we can write the right-hand side of (13.11) as

\int

B

E[1A]dP = E [1A]

\int

B

dP= P(A)P(B).

Therefore, (13.11) holds forX= 1 A and anyB\in \mathcal{G}.

We can extend this to simple functions, nonnegative functions, and finally to

real-valued functions. \hfill $\square$

\textbf{Example 13.3.1 (Iterated Expectation)} \\

For any random variables X and Y , apply Theorem 13.8 with \mathcal{H}={ \emptyset ,\Omega} and \mathcal{G}= \sigma(Y) S i n c e

E[Z\vert\mathcal{H}] is a constant and is equal toE[Z] for any random variable Z ,we getE[E[X\vertY]] =E [X].
